\chapter{总结与展望}
寻找新物理和理解QCD非微扰是目前粒子物理领域内的核心任务。而粲物理，特别是粲味强子弱衰变对于这两个核心任务起到关键作用：一方面，此类过程作为弱衰变，对可能存在的高能标新物理非常敏感，另一方面，过程中初末态强子化由非微扰 QCD 所主导，是探究 QCD的重要平台。粲味强子弱过程分为遍举衰变和单举衰变，后者理论研究方法比较成熟，可以从第一性原理出发系统的求解幂次修正的形式，因而其理论误差更加可控。相比于B介子和底重子，由于粲夸克质量较小，按QCD耦合常数展开的微扰修正以及按重夸克质量负幂次展开的幂次修正都更加显著，因此粲味强子弱衰变为研究高阶修正提供了更加理想的环境。因此粲味强子单举衰变的研究具有其独特意义。

与此同时，粒子物理领域内长期存在一些疑难与反常，并引起领域内的广泛关注。对于更清楚的认识甚至解决这些疑难，粲味强子单举衰变的研究可以起到关键性的作用。其中之一是粲强子寿命的理论预言与实验测量不一致，原因在于寿命理论计算依赖非微扰参数，目前理论工作中的非微扰参数均有较大的模型依赖。因此，基于粲味强子举衰变过程的观测量和对应实验数据，实现对非微扰矩阵元模型无关的抽取，是理解甚至解决粲强子寿命疑难的关键所在。

本文在以下三个方面对理解粲强子寿命疑难进行了努力：

首先理论上，基于重夸克有效理论和算符乘积展开，讨论高阶修正对非微扰强子矩阵元的影响。高阶修正主要包括重夸克质量负幂次展开的幂次修正和QCD耦合常数展开的微扰修正。关于重夸克质量负幂次展开的幂次修正，介绍了$\mrm{B\to X_{c}e\nu}$过程树图水平量纲三和量纲五算符对应结构函数的匹配与相空间解析积分，得到对应的电子能谱解析形式。在此基础上，通过取末态粲夸克质量趋于零的极限，可得$\mrm{B\to X_{u}e\nu}$过程树图水平的电子能谱解析形式。关于QCD耦合常数展开的微扰修正，在单圈水平上，本文介绍了$\mrm{B\to X_{c}e\nu}$过程的结构函数，并通过取末态粲夸克质量趋于0的极限，得到$\mrm{B\to X_{u}e\nu}$过程的结构函数，在此基础上，对相空间进行解析积分，得到对应的电子能谱解析形式。将上述的电子能谱转换到粲夸克半轻单举衰变过程中，即得到$\mrm{D\to X e^{+}\nu}$过程电子能谱的解析形式。由于重夸克有效理论的局限性，需要对电子能谱进行相空间加权积分进行“平滑”处理后，方可与实验观测量进行对比。本文讨论了$\mrm{D\to X_s e^{+}\nu}$和$\mrm{D\to X_d e^{+}\nu}$两类过程，粲强子半轻单举衰变的末态夸克包括奇异夸克，其质量接近非微扰能标，因此需要作为软自由度进行保留，对于$\mrm{D\to X_s e^{+}\nu}$，末态奇异夸克的质量效应被考虑在内。基于$\mrm{D\to X e^{+}\nu}$过程的电子能谱，通过对其做相空间加权解析积分，给出关于衰变宽度、电子能量原点矩的理论表达式，它们依赖于非微扰强子矩阵元。

实验上，美国的CLEO合作组2009年测量了实验室系下三种D介子半轻单举衰变过程的电子能谱，我国的BESIII合作组于2018年更新了$\mrm{D_s^+}$介子半轻单举衰变过程的电子能谱。本文从粲介子半轻单举衰变的电子能谱出发，利用重夸克有效理论的预言，在红外区域对其做拟合外推；关于系统误差，对衰变分支比的系统误差进行延拓并假设每个bin的系统误差全关联，最终得到实验室系下完整的电子能谱。基于此，对其做关于电子能量的加权求和得到电子能量矩和衰变宽度的实验测量值。除此之外，还讨论了电子能量矩之间的关联系数，最终选择衰变宽度，电子能量一阶原点矩，电子能量第二、三和四阶中心矩作为数据点与对应的理论表达式一齐进行全局拟合。

唯象学分析上，由于衰变宽度对粲夸克质量高度敏感，因此本文基于三种粲夸克质量方案，给出对应方案下粲强子半轻单举衰变宽度、电子能量一阶原点矩电子能量第二、三和四阶中心矩的理论表达式，与对应实验观测值一齐进行全局拟合；关于幂次修正对于非微扰参数拟合抽取带来的误差贡献，本文分为两种情形进行全局拟合，分别是精确至量纲五算符和精确至量纲六算符贡献。关于理论输入参数的误差，本文仅考虑由于重整化能标跑动引起的强相互作用耦合常数的误差，通过2000次MC撒点进行全局拟合，取其非微扰参数估计值分布的标准误差作为理论输入参数$\alpha_s$的误差对非微扰参数估计值的误差贡献。

唯象学分析结果表明：
\begin{itemize}
    \item  对于三种粲夸克质量方案全局拟合，引入量纲六算符非微扰矩阵元贡献之后，平均最小卡方值迅速降低，表明量纲六算符非微扰矩阵元在衰变宽度的理论计算中起到关键性的作用。
    \item  在1S质量方案下，平均最小卡方值小于1，表明该质量方案下的拟合具有良好的效果，且1S质量方案和$\overline{\mrm{MS}}$方案下非微扰参数的拟合抽取值保持一致，最终本文推荐1S质量方案下非微扰参数的估计值作为推荐。
    \item  在1S质量方案下，基于非微扰参数的拟合估计值，本文计算了D介子半轻单举衰变宽度的微扰修正和幂次修正，呈现出极好的收敛性。
    \item  基于本文关于非微扰参数的拟合估计值，本文对粲强子衰变宽度进行了定性分析，结果表明使用本文关于非微扰参数的估计值，三种D介子寿命均更接近实验观测量值。
\end{itemize}

本文首次对粲强子半轻单举衰变过程中的量纲五算符非微扰强子矩阵元实现模型无关的高精度拟合抽取：该拟合结果可以直接用于粲强子寿命的理论计算；提出1S质量方案是目前计算粲强子寿命的最适合的质量方案：我们在1S方案方案下给出粲强子半轻单举衰变宽度精确至三圈图微扰修正的数值结果，展示出极好的收敛性。

本文工作还存在很多地方需要提升，主要包括：
\begin{itemize}
    \item 理论上：量纲五算符的单圈修正贡献不可忽略，由于逐阶计算的难度和计算量将爆炸式地增长，因此本文并未讨论该修正项对非微扰参数估计值的影响，但未来考虑量纲五算符的单圈修正是必要的。
    \item 数据方面：一方面，从实验数据中抽取衰变宽度、电子能量矩的工作需要实验粒子物理学家进行更专业的分析；另一方面，实验上测量D介子质心系下的电子能谱是必要的，这将有助于对相空间做切割，使得构建更多的数据点成为可能。
    \item 唯象学分析上，随着未来理论和实验精度的不断提升，更加系统全面的拟合手段的引入是必要的，比如贝叶斯分析。
\end{itemize}

对于该系列工作，从实验角度来看，BESIII 和即将问世的超级陶粲工厂预计将提供更精确的总衰变率和微分衰变率的测量，这对于更好地确定重夸克有效理论参数至关重要。此外，我们建议实验与以往的研究相比，纳入以下改进：
\begin{itemize}
\item 在 D介子的质心系而不是实验室系中测量电子能谱：这使得理论物理学家可以直接利用结果，而无需使用洛伦兹变换，洛伦兹变换仅对整个相空间中的电子能量矩有效。因此，可以得到具有更高能量截断的电子能量矩，为重夸克有效理论参数确定提供更多的实验可观测量。
\item 直接测量电子能量矩：从分bin的电子能量谱中获得电子能量矩需要假设每个箱内的特定分布，引入不必要的不确定性。
\item 利用单举过程强子系统 $X$ 的信息，并重建轻子对 $\left(q^2\right)$ 和强子系统 $\left(M_X\right)$ 的不变质量平方。特别是 $q^2$ 矩的结果，将增强具有重参数化不变性的HQET算符的非微扰强子矩阵元的确定。通过包括足够多的可观测量，如具有各种截断的 $E_e, ~ q^2$ 和 $M_X$ 矩，就可以确定量纲七算符的非微扰强子矩阵元。
\item 对于$\mrm{D\to \mu X}$的衰变道进行测量，所有可观测量的 $\mu/e$比值可以显著减少理论不确定性，使它们成为标准模型的精确检验可观测量。
\item 区分$\mrm{D\to X_{s} e^+\nu}$ 中的$\mrm{D\to X_{d} e^+\nu}$,这将使得抽取单举版本的CKM矩阵元$\mrm{V_{cs}}$和$\mrm{V_{cd}}$成为可能，并与遍举衰变中抽取出的CKM矩阵元进行比较，以检验CKM机制。
\end{itemize}
最后，该工作可以推广至味道改变中性流$\mrm{c\to u ll}$过程，检验轻子味道普世性的破坏以及角分布反常，这将有助于理解B介子衰变中相应的反常现象。